\documentclass[a4paper,12pt]{article}
\usepackage{listings}
\usepackage{xcolor}
\usepackage{geometry}
\usepackage{amsmath}
\usepackage{booktabs}
\usepackage{graphicx}

\geometry{margin=1in}

\definecolor{codegray}{gray}{0.9}
\lstset{
  backgroundcolor=\color{codegray},
  basicstyle=\ttfamily\small,
  frame=single,
  breaklines=true,
  keywordstyle=\color{blue},
  commentstyle=\color{green!50!black},
  numbers=left,
  numberstyle=\tiny,
  stepnumber=1,
  numbersep=8pt,
  tabsize=2,
  showspaces=false,
  showstringspaces=false,
  captionpos=b
}

\begin{document}

\begin{center}
{\LARGE \textbf{Experiment: K-Map Based Logic Gate Design using VAMAN Board}} \\[6pt]
\textbf{Course: Digital Logic Design using Verilog HDL}
\end{center}

\section*{Objective}
To implement a simplified Boolean function derived from its Karnaugh Map (K-Map) and verify it on the VAMAN FPGA development board.

\section*{Apparatus Required}
\begin{itemize}
  \item VAMAN FPGA Development Board (Spartan-6)
  \item USB JTAG Programmer
  \item On-board switches and LEDs
  \item Xilinx ISE / Vivado Software
\end{itemize}

\section*{Boolean Function}
The logic expression obtained from K-map simplification is:
\[
f = (b + c)(a + \overline{c})
\]
This expression represents the minimized Product of Sums (PoS) form.

\section*{Truth Table}
\begin{center}
\begin{tabular}{ccc|c}
\toprule
A & B & C & F \\ 
\midrule
0 & 0 & 0 & 0 \\
0 & 0 & 1 & 0 \\
0 & 1 & 0 & 1 \\
0 & 1 & 1 & 0 \\
1 & 0 & 0 & 0 \\
1 & 0 & 1 & 1 \\
1 & 1 & 0 & 1 \\
1 & 1 & 1 & 0 \\
\bottomrule
\end{tabular}
\end{center}

\section*{K-Map Simplification}
\[
\begin{array}{c|cccc}
AB \backslash C & 0 & 1 \\ \hline
00 & 0 & 0 \\
01 & 1 & 0 \\
11 & 1 & 0 \\
10 & 0 & 1 \\
\end{array}
\]
The grouped 1s in the K-map lead to:
\[
f = (b + c)(a + \overline{c})
\]

\section*{Verilog Code}
\begin{lstlisting}[language=Verilog, caption={Verilog Code for K-Map Simplified Logic Function}]
module kmap_gate_vaman(
    input wire a,
    input wire b,
    input wire c,
    output wire f
);

assign f = (b | c) & (a | ~c);

endmodule
\end{lstlisting}

\section*{Pin Configuration (VAMAN Board Example)}
\begin{center}
\begin{tabular}{lll}
\toprule
\textbf{Signal} & \textbf{FPGA Pin} & \textbf{Description} \\
\midrule
a & SW0 & Input Switch 0 \\
b & SW1 & Input Switch 1 \\
c & SW2 & Input Switch 2 \\
f & LED0 & Output LED 0 \\
\bottomrule
\end{tabular}
\end{center}

\section*{Result}
The output LED glows whenever the logic expression \((b + c)(a + \overline{c})\) evaluates to HIGH, confirming correct implementation of the simplified K-map expression.

\section*{Conclusion}
The Verilog module successfully implements the minimized Boolean function obtained from the K-map.  
This verifies the use of logical reduction techniques for efficient hardware implementation on the VAMAN FPGA board.

\end{document}
